
Modern neural networks, including large language models, are poorly calibrated: their confidence scores do not reliably reflect accuracy \citep{guo2017calibration}. On factual tasks such as TruthfulQA \citep{lin2022truthfulqa}, models express high confidence even when generating false statements. Post-hoc calibration methods, particularly temperature scaling, can reduce Expected Calibration Error by adjusting the softmax temperature to better align confidence with accuracy.

TruthfulQA organizes 817 questions across 38 semantic categories ranging from Finance and Economics to Misconceptions and Myths. A natural hypothesis emerges: if models exhibit different confidence-accuracy relationships across these domains, category-specific calibration should outperform a single global temperature. This intuition aligns with findings that LLM uncertainty varies by task type \citep{xiong2023uncertainty}.

Despite extensive work on neural network calibration and growing interest in LLM uncertainty quantification, no prior study has systematically tested whether category-level calibration structure exists in truthfulness benchmarks or whether it can be exploited via per-category temperature scaling.

This paper addresses this gap through experiments testing cluster-specific temperature scaling on TruthfulQA using Llama-2-7B. The 38 categories are grouped into 7 semantic clusters, and three linked hypotheses are tested:

\begin{enumerate}
\item \textbf{Existence (h-e1):} Does calibration error vary significantly across clusters?
\item \textbf{Mechanism (h-m1):} Do different clusters produce distinct confidence distributions?
\item \textbf{Mechanism (h-m2):} Do different distributions require different temperature parameters?
\end{enumerate}

The results confirm the first two hypotheses while falsifying the third. Significant ECE variation across clusters is observed (F=8.45, p=0.00012), yet when optimizing temperature per cluster, all seven clusters converge to the same value: T=10.0, the upper bound of the search range.
